%!TEX program = xelatex
%!TEX encoding = UTF-8
%Definições de Ano, Semestre, número e data da avaliação
\def\numeroteste{2}
\def\modulo{1}
\def\dataavaliacao{12/04/2024}

\documentclass[12pt]{exam}

\usepackage{caption}
\usepackage{amssymb}
\usepackage{amsmath,amsfonts,amsthm,amstext}
\usepackage[brazil]{babel}
% \usepackage[latin1]{inputenc}
\usepackage{graphicx}
\graphicspath{{/home/jfreitas/GitHub_Repos/Algebra-1/Pictures/}{D:/OneDrive - unb.br/imagens-latex/}}
\usepackage{enumitem}
\usepackage{multicol}
\usepackage{answers}
\usepackage{tikz,ifthen}
\usetikzlibrary{lindenmayersystems}
\usetikzlibrary[shadings]
\def\ano{2024}
\def\semestre{1}
\def\disciplina{Álgebra 1}
\def\nomeabreviado{Álgebra 1}
\def\turma{1}

\newcommand{\im}{{\rm Im\,}}
\newcommand{\dlim}[2]{\displaystyle\lim_{#1\rightarrow #2}}
\newcommand{\minf}{+\infty}
\newcommand{\ninf}{-\infty}
\newcommand{\cp}[1]{\mathbb{#1}}
\newcommand{\sub}{\subseteq}
\newcommand{\n}{\mathbb{N}}
\newcommand{\z}{\mathbb{Z}}
\newcommand{\rac}{\mathbb{Q}}
\newcommand{\real}{\mathbb{R}}
\newcommand{\complex}{\mathbb{C}}

\newcommand{\vesp}[1]{\vspace{ #1  cm}}

\newcommand{\compcent}[1]{\vcenter{\hbox{$#1\circ$}}}
\newcommand{\comp}{\mathbin{\mathchoice
        {\compcent\scriptstyle}{\compcent\scriptstyle}
        {\compcent\scriptscriptstyle}{\compcent\scriptscriptstyle}}}
\renewcommand{\sin}{{\rm sen\,}}
\renewcommand{\tan}{{\rm tg\,}}
\renewcommand{\csc}{{\rm cossec\,}}
\renewcommand{\cot}{{\rm cotg\,}}
\renewcommand{\sinh}{{\rm senh\,}}
\newcounter{exercicios}
\setcounter{exercicios}{0}
\newcommand{\questao}{
    \addtocounter{exercicios}{1}
    \noindent{\bf Questão \arabic{exercicios}: }}

\newcommand{\solucao}{
    \noindent{\bf Solução: }}

\newcommand{\resp}[1]{
    \noindent{\bf Exerc{\'\i}cio #1: }}

\extrafootheight[.25in]{.5in}
\footrule
\lfoot{Teste \numeroteste\ - Módulo \modulo\ - \nomeabreviado\ - Turma \turma\ - \semestre$^o$/\ano}
\cfoot{}
\rfoot{P\'agina \thepage\ de \numpages}

\begin{document}

    \begin{figure}[h]
        \begin{minipage}[c]{1.7cm}
            \includegraphics[width=1.7cm]{unb.pdf}
        \end{minipage}
        \hspace{0pt}
        \begin{minipage}[c]{4in}
            {Universidade de Brasília} \\
            {Departamento de Matemática}
        \end{minipage}
    \end{figure}
    \vesp{-1}
    \hrule
    \begin{center}
        {\Large\bf \disciplina\ - Turma \turma}  \\
        \vesp{0.2} {\large\bf Teste \numeroteste\ - Módulo\ \modulo\ -\ \dataavaliacao}
    \end{center}

    \noindent{\bf \underline{Nome: \hspace{9.8cm} Mat.:\qquad \
            \hspace{2.1cm}}}\\
    \vspace*{.01cm}

    \noindent{\bf \underline{Nome: \hspace{9.8cm} Mat.:\qquad \
            \hspace{2.1cm}}}
    \vspace{.4cm}

    \questao{} Em $\z$ defina a relação $R$ por
    \[
    a R b \mbox{ se, e somente se, } |a + b| \mbox{ é par}
    \]
    para $a$, $b \in \z$.
    \begin{enumerate}[label={\alph*})]
        \item Mostre que $R$ é uma relação de equivalência em $\z$.

        \item Determine todas as classes de equivalência dessa relação.
    \end{enumerate}
    \solucao
    \begin{enumerate}[label={\alph*})]
        \item Seja $a \in \z$. Temos
        \[
            |a + a| = |2a| = 2|a|
        \]
        que é sempre um número par. Logo $aRa$.

        Agora suponha que $aRb$. Assim
        \[
            |a + b| = 2k
        \]
        para algum $k \in \z$. Com isso temos
        \[
            |b + a| = |a + b| = 2k.
        \]
        Logo $bRa$.

        Finalmente, suponha que $aRb$ e $bRc$. Assim
        \begin{align}
            |a + b| = 2k\label{primeiraequacao}\\
            |b + c| = 2l\label{segundaequacao}
        \end{align}
        para $k$, $l \in \z$.

        De \eqref{primeiraequacao} segue que se $a$ é par, então $b$ é par. Também obtemos que se $a$ é ímpar, então $b$ é impar.

        Agora, de \eqref{segundaequacao} segue que se $b$ é par, então $c$ é par. Também obtemos que se $b$ é ímpar, então $c$ é impar.

        Assim temos que se $a$ é par, então $b$ é par e com isso $c$ também deve ser par. Logo $|a + c|$ será par. Por outro lado, se $a$ é ímpar, então $b$ é ímpar e com isso $c$ também será ímpar. Logo $|a + c|$ será par. Independente do caso, sempre teremos $|a + b|$ par e com isso $aRc$.

        Portanto $R$ é uma relação de equivalência sobre $\z$.

        \item As classes de equivalência de $R$ serão:
        \begin{align*}
            \overline{0} &= \{x \in \z : |x + 0| = 2k,\ k \in \z\} = \{x \in \z : |x| = 2k, \ k \in \z\} = \{0, \pm 2, \pm 4, \pm 6, \dots\}\\
            \overline{1} &= \{x \in \z : |x + 1| = 2k,\ l \in \z\} = \{x \in \z : |x + 1| = 2l, \ k \in \z\} = \{\pm 1, \pm 3, \pm 5, \dots\}
        \end{align*}
        Portanto existem somente duas classes de equivalência.
    \end{enumerate}
    \vspace*{.75cm}

    \questao{} Considere a seguinte relação sobre $\z$:
    \[
    x \sim y \mbox{ quando }  x^2 - y^2 = 4k, \mbox{ para algum } k \in \z.
    \]

    \begin{enumerate}[label={\alph*})]
        \item Mostre que $\sim$ é relação de equivalência.

        \item Determine todas as classes de equivalência dessa relação.
    \end{enumerate}

        \solucao
    \begin{enumerate}[label={\alph*})]
        \item Seja $x \in \z$. Temos
        \[
        x^2 - x^2 = 0 = 4\cdot 0
        \]
        logo $x \sim x$.

        Agora, suponha que $x \sim y$. Daí
        \[
        x^2 - y^2 = 4k
        \]
        para algum $k \in \z$. Assim
        \[
        y^2 - x^2 = -(x^2 - y^2) = -4k = 4(-k),
        \]
        logo $y \sim x$.

        Finalmente, suponha que $x \sim y$ e que $y \sim z$. Assim
        \begin{align}
            x^2 - y^2 = 4k\label{equacao1}\\
            y^2 - z^2 = 4l\label{equacao2}
        \end{align}
        para $k$, $l \in \z$. Somando \eqref{equacao1} e \eqref{equacao2} obtemos
        \[
        x^2 - z^2 = (x^2 - y^2) + (y^2 - z^2) = 4k + 4l = 4(k + l).
        \]
        Logo $x \sim z$.

        Portanto $\sim$ é uma relação de equivalência sobre $\z$.

        \item As classes de equivalência de $\sim$ são:
        \begin{align*}
            \overline{0} &= \{x \in \z \mid x^2 - 0^2 = 4k,\ k \in \z\} = \{x \in \z \mid x^2 = 4k,\ k \in \z\} = \{0, \pm 2, \pm 4, \pm 6, \dots\}\\
            \overline{1} &= \{x \in \z \mid x^2 - 1^2 = 4l,\ l \in \z\} = \{x \in \z \mid x^2 = 4l + 1,\ l \in \z\} = \{\pm 1, \pm 3, \pm 5, \dots\}.
        \end{align*}
    \end{enumerate}

    Portanto existem somente duas classes de equivalência.

\end{document}
